\section{Discusión}
\label{sec:discussion}

Las decisiones de diseño en Frankenstein Transformer reflejan varias tensiones de ingeniería e investigación en las herramientas modernas de aprendizaje profundo.

La campaña reportada en la Sección~\ref{sec:empirical-evaluation} convierte varias de esas tensiones en mediciones en lugar de opiniones. El esquema estricto rindió operativamente: 868 configuraciones se entrenaron, desplegaron y evaluaron a través del mismo contrato, y los fallos que quedaron se concentraron en cuatro mezcladores y en la serialización de checkpoints, no en la deriva de configuración. El sistema de optimizadores prefijados resultó expresivo pero, a 100 épocas, indeciso---una docena de optimizadores convergió a la misma pérdida---lo que es en sí mismo un argumento a favor del esquema: solo una barrida controlada puede mostrar que una dimensión de diseño no importa a una escala dada.

\subsection{Compromisos de Diseño Basado en Esquemas}

El enfoque basado en esquemas proporciona beneficios significativos de reproducibilidad al imponer contratos explícitos y fallar rápidamente ante configuraciones inválidas. Sin embargo, este enfoque también introduce rigidez: agregar nuevas arquitecturas u optimizadores requiere extensiones del esquema en lugar de argumentos de línea de comandos flexibles. El sistema de hiperparámetros prefijados permite un control detallado pero aumenta la complejidad de configuración para usuarios acostumbrados a interfaces más simples.

La decisión de imponer \texttt{additionalProperties: false} en todos los niveles del esquema elimina la absorción silenciosa de parámetros que ha afectado a sistemas de configuración anteriores, pero esta rigurosidad requiere un mantenimiento cuidadoso del esquema al extender las capacidades del sistema. Cada nuevo mecanismo de atención o variante de optimizador debe integrarse adecuadamente en el marco de validación, incluyendo definiciones de campos del esquema con tipos y restricciones apropiados, mapeo de hiperparámetros prefijados para grupos específicos de optimizadores, valores predeterminados alineados con las mejores prácticas de investigación y cadenas de documentación para la representación en la interfaz web.

\subsection{Cobertura Arquitectónica y Vacíos}

Las cincuenta y cuatro arquitecturas de mezclador implementadas abarcan las principales direcciones de investigación en modelado de secuencias, pero ciertos vacíos permanecen. El sistema carece de arquitecturas híbridas recientes como Griffin \citep{soares_griffin_2024} y Jamba \citep{gu_jamba_2024}, que combinan compuertas con modelos de espacio de estados. El enrutamiento MoE (Mezcla de Expertos) está implementado para capas FFN pero no para el cómputo de atención, donde trabajos recientes han mostrado beneficios \citep{lewis_moe_2021}.

La cobertura de atención dispersa es amplia. Los dos métodos en tiempo de inferencia (FASA, SpargeAttn) conservan su carácter sin entrenamiento en el sentido de que su regla de dispersidad nunca se aprende; en cambio, la implementación entrena la capa como atención estándar densa sobre las mismas proyecciones y activa la ruta de selección/poda dispersa solo en evaluación, de modo que ambos mixers pueden entrenarse desde cero y su comportamiento disperso es exactamente el descrito en los artículos originales.

La cobertura de mecanismos de compuerta y peso rápido abarca las categorías principales, incluidas las seis variantes Falcon \citep{zhang_falcon_2026}. El macro-diseño del flujo residual está cubierto por Attention Residuals \citep{sun_attnres_2026} y por el flujo residual de $n$ streams restringido por variedad (mHC) \citep{xie_mhc_2025}, cuya proyección Sinkhorn--Knopp sobre el politopo de Birkhoff restaura la propiedad de mapeo de identidad que pierden las hiper-conexiones sin restricción; mHC está validado solo al nivel de los experimentos MoE publicados de 3B--27B y se beneficiaría de reproducciones independientes a escalas menores (Apéndice~\ref{sec:annex-mhc}).

\subsection{Fragmentación del Panorama de Optimizadores}

El soporte para treinta y ocho familias de optimizadores en ocho categorías demuestra exhaustividad pero también resalta el estado fragmentado de la investigación en optimización. Los usuarios enfrentan una complejidad significativa de decisión al elegir entre métodos de reducción de varianza (Adan, MARS), variantes eficientes en memoria (GaLore, Adafactor, APOLLO) y enfoques conscientes de curvatura (Sophia, Shampoo). El sistema de hiperparámetros prefijados, aunque poderoso, requiere comprender qué parámetros son relevantes para cada clase de optimizador.

La calidad de implementación varía entre optimizadores: los métodos clásicos (AdamW, SGD con momento) están altamente optimizados en PyTorch, mientras que los métodos más nuevos (Muon, Turbo-Muon, AF-Muon, SOAP) pueden requerir implementaciones personalizadas que afectan la estabilidad numérica y las características de rendimiento.

\subsection{Consideraciones de Despliegue y Producción}

El pipeline de cuantización demuestra preocupaciones prácticas de despliegue pero realiza compromisos de ingeniería específicos. El empaquetado ternario de pesos reduce el almacenamiento a aproximadamente $1.58$ bits por parámetro, pero esta compresión agresiva puede degradar el rendimiento, especialmente para modelos más pequeños donde el error de cuantización es más significativo. El exportador safetensors ofrece además formatos afines por bloques \texttt{Q8\_0} y \texttt{Q4\_0} (8.5 y 4.5 bits por peso) y aplica los tres solo a pesos lineales bidimensionales cuya dimensión de entrada es divisible por 32; los embeddings, los parámetros de normalización y los sesgos permanecen en FP32, por lo que la compresión efectiva depende de la fracción lineal del modelo. El runtime packed-linear mantiene el artefacto compacto a costa de un paso de des-cuantización en cada pase forward.

Los flujos de trabajo SBERT proporcionan utilidad práctica para tareas de similitud semántica y recuperación, pero la implementación asume estrategias de pooling estándar (token CLS, pooling promedio). Avances recientes como los embeddings Matryoshka \citep{muennighoff_matryoshka_2021} y los refinamientos de aprendizaje contrastivo \citep{chen_supervised_2020} aún no están incorporados.

\subsection{Desafíos de Integración y Extensibilidad}

La extensibilidad se gestiona mediante un registro declarativo de componentes en lugar de un despacho ad-hoc. Cada componente conectable---mezclador de secuencia, optimizador, normalización, activación, estrategia residual, codificación posicional y embedding de token---se describe con un \texttt{ComponentSpec} inmutable registrado a través de \path{src/registry/}, y las enumeraciones del esquema, los fragmentos de parámetros por componente, la referencia de prefijos del optimizador y las reglas condicionales de prefijo se ensamblan desde ese registro antes de la validación (Sección~\ref{sec:system-design}). Un plugin de terceros registra el mismo tipo de spec mediante la API pública de autoría, ya sea desde la lista \texttt{plugins:} de la configuración o a través del grupo de entry-points \texttt{frankenstein.plugins}, de modo que valida estrictamente, aparece en el constructor generado desde el esquema y es despachado por el mismo código central que los integrados. La presión de mantenimiento es, por tanto, el tamaño del catálogo de componentes---cuarenta y nueve mezcladores y treinta y ocho optimizadores---y no la lógica de despacho, que no contiene ramas por optimizador ni una escalera de selección por mezclador.

La interfaz de configuración web proporciona mejoras significativas de usabilidad pero introduce complejidad de despliegue: ejecutar Streamlit junto con trabajos de entrenamiento requiere recursos adicionales y consideraciones de infraestructura que pueden no ser apropiadas para todos los entornos, particularmente clústeres HPC sin acceso web.

